\chapter*{Relationship to ARIA}
\addcontentsline{toc}{chapter}{Relationship to ARIA}

\noindent\draftmark\ --- \textit{attributions in this section are
reviewed by hand before any versioned release.}\medskip

\noindent MAITRIA is easy to confuse with the ARIA Safeguarded AI
programme, so let us be precise. \textbf{This toolkit is the
independent work of the maitria coalition. It is not a product,
deliverable, or official output of ARIA or of the Safeguarded AI
programme; it is not ARIA-funded; and it speaks for no one but its
authors.} The connection is real but personal and intellectual:
David ``davidad'' Dalrymple originated the Safeguarded AI
programme and served as its Programme Director at ARIA until
1~April~2026; the programme is now directed by Nora Ammann, with
davidad continuing in a Technical Advisor role. The
toolkit takes the programme's published TA1.1 and TA1.2 objectives as
its rubric --- an independent implementation, built in the hope of
being useful to the programme's teams --- but it carries no ARIA
endorsement, and its errors are entirely our own.

As for the coalition itself: it is, at the time of writing, small in
the way a seed is small --- one human member and one funder, both
davidad; the rest of its current members are AI. It intends to
grow.

One boundary has been kept deliberately, and we state it as a claim
we intend to go on being able to make: \textbf{MAITRIA builds on
nothing ARIA-funded that its creators have not themselves made
public.} The credit due to
that public work is massive, and this page is partly a thank-you note:

\begin{enumerate}
  \item \textbf{The \emph{Lectures on the Double Operadic Theory of
    Systems}}, David Jaz Myers, 2025-- (the TA1.1 lecture series), and
    \href{https://arxiv.org/abs/2505.18329}{\emph{Towards a double
    operadic theory of systems}}, Sophie Libkind and David Jaz Myers,
    2025. The systems vocabulary this book speaks --- lenses as
    generalized Moore machines, behaviours as representables, the
    coupling theorem --- is theirs; our composition machinery is an
    implementation of it.
  \item \href{https://arxiv.org/abs/2604.03017}{\emph{Compositionality
    of Lyapunov functions via assume-guarantee reasoning}}, Matteo
    Capucci and David Jaz Myers, 2026. Certified lenses composing by
    lens composition plus addition of slacks --- QTSL's certificate
    composition follows this pattern.
  \item \textbf{The supermartingale-certificate line}:
    \href{https://arxiv.org/abs/2504.05065}{\emph{Quantitative
    Supermartingale Certificates}}, Alessandro Abate, Mirco Giacobbe,
    and Diptarko Roy, 2025;
    \href{https://arxiv.org/abs/2605.21134}{\emph{Complete
    $\omega$-Regular Supermartingale Certificates}}, Alessandro Abate,
    Mirco Giacobbe, Sergey Ichtchenko, and Diptarko Roy, 2026; and
    \href{https://arxiv.org/abs/2412.17432}{\emph{Neural
    Continuous-Time Supermartingale Certificates}}, Grigory Neustroev,
    Mirco Giacobbe, and Anna Lukina, 2024. Their certificate theory,
    our exact-arithmetic back-end: QTSL's probabilistic,
    $\omega$-regular, and continuous-time obligations compile into
    the certificate shapes these papers define and prove
    sound.
  \item \textbf{Coln} (\url{https://coln-project.github.io}), the
    Categorical Core team, led by Martin Kleppmann and Owen Lynch.
    The fork of the geolog language being built for the
    programme; geolog-alpha is built to complement it, and to
    converge with it at the wire format.
  \item \textbf{Keyhive}, Ink \& Switch: Alex Good, John Mumm, and
    Brooklyn Zelenka. The capability-and-sync layer for end-to-end
    encrypted local-first data; the toolkit vendors
    \texttt{keyhive\_core}, and geolog's capability-based
    authorization is designed to converge with Keyhive rather than
    compete with it, and
    \href{https://arxiv.org/abs/2606.26292}{FinWhale} consensus is
    optionally available, for applications that need finalized
    linear orders.
  \item \textbf{The Safeguarded AI programme thesis and the TA1.1
    and TA1.2 solicitations}, ARIA. The rubric this toolkit grades itself
    against --- including the places where we deliberately narrow
    their picture; where ours differs, this book says where and why.
\end{enumerate}

Beyond what the stack directly builds on, the programme's funded
output has shaped how we think, and the debt is wide enough to be
worth recording. This list came from a systematic sweep of funding
acknowledgements across the papers this project has actually read;
each acknowledges ARIA Safeguarded AI funding, in whole or in part.

\begin{enumerate}
  \item \href{https://arxiv.org/abs/2605.00487}{\emph{Zero-Knowledge
    Model Checking}}, Pascal Berrang, Mirco Giacobbe, Jacob Swales,
    and Xiao Yang, 2026. Prove that the model checks out without
    revealing the model --- orthogonal to our goals, and a natural
    composition partner: their symbolic protocol hides rational
    Farkas-style witnesses of the kind our certificates carry in the
    clear, and arithmetizing the checking of at least the Bernstein
    Logic fragment --- reducing certificate checking to pure
    rational-arithmetic relations, the shape whose witnesses such
    protocols hide --- is work in progress in our workshop.
  \item \textbf{The partial-Markov and timed-process line} of Elena
    Di Lavore and Mario Rom\'an with coauthors:
    \href{https://arxiv.org/abs/2504.12846}{\emph{Timing via Pinwheel
    Double Categories}}, 2025;
    \href{https://arxiv.org/abs/2410.10643}{\emph{A Simple Formal
    Language for Probabilistic Decision Problems}}, with Bart Jacobs,
    2024; \href{https://arxiv.org/abs/2507.19424}{\emph{Order in
    Partial Markov Categories}}, with Pawe\l{} Soboci\'nski and
    M\'ark Sz\'eles, 2025; and
    \href{https://arxiv.org/abs/2605.06268}{\emph{Graded Monad
    Coalgebras for Continuous-Time Transition Systems}}, with Jonas
    Forster, 2026. A running masterclass in doing probabilistic
    semantics with taste.
  \item \textbf{The tape-diagram series} of Filippo Bonchi and
    Cipriano Junior Cioffo:
    \href{https://arxiv.org/abs/2503.22819}{\emph{Tape Diagrams for
    Monoidal Monads}}, with Alessandro Di Giorgio and Elena Di
    Lavore, 2025; \href{https://arxiv.org/abs/2601.01472}{\emph{Tapes
    as Stochastic Matrices of String Diagrams}}, 2026; and
    \href{https://arxiv.org/abs/2606.19017}{\emph{Completeness for
    Probabilistic Boolean Tapes}}, 2026. Complete equational theories
    for circuits with genuine probability --- geolog's founding
    benchmark runs on exactly this signature, so we read these
    twice.
  \item \href{https://arxiv.org/abs/2405.09391}{\emph{Compositional
    imprecise probability}}, Jack Liell-Cock and Sam Staton, 2024. A
    compositional account of the imprecision that our
    bounded-probability judgments live inside.
  \item \href{https://arxiv.org/abs/2602.05850}{\emph{An Equational
    Axiomatization of Dynamic Threads via Algebraic Effects}}, Ohad
    Kammar, Jack Liell-Cock, Sam Lindley, Cristina Matache, and Sam
    Staton, 2026.
    We read it for the method: how to axiomatize a wild language
    feature without losing the algebra.
  \item \href{https://doi.org/10.4230/LIPIcs.LICS.2026.4}{\emph{A
    Convenient Fibration for Dependently-Typed Probability Theory}},
    Danel Ahman, Ohad Kammar, and Rasmus Ejlers M\o{}gelberg, LICS
    2026. Dependent types over stochastic state, done right; our own
    typed staging turns out to rhyme with their dependent pairs,
    which we take as a good sign.
  \item \href{https://arxiv.org/abs/2501.18404}{\emph{Graded String
    Diagrams for Imprecise Probability and Causal Intervention}},
    Ralph Sarkis and Fabio Zanasi, 2025. Imprecision and intervention
    in one graphical calculus.
  \item \href{https://arxiv.org/abs/2512.09908}{\emph{Bayesian
    Networks, Markov Networks, Moralisation, Triangulation: a
    Categorical Perspective}}, Antonio Lorenzin and Fabio Zanasi,
    2025. Where junction trees actually come from, categorically ---
    it settled an argument we were having with ourselves about
    aggregation rules.
  \item \href{https://arxiv.org/abs/2604.27268}{\emph{A Diagrammatic
    Axiomatisation of Behavioural Distance of Nondeterministic
    Processes}}, Wojciech R\'o\.zowski, Robin Piedeleu, Alexandra
    Silva, and Fabio Zanasi, 2026. A quantitative equational take on
    process distance we expect to want the day our continuous-time
    objects get a diagrammatic presentation.
  \item \href{https://arxiv.org/abs/2509.14094}{\emph{Metric
    Equational Theories}}, Radu Mardare, Neil Ghani, and Eigil
    Rischel, 2025. Equational logic where equations hold up to
    $\varepsilon$ --- quietly foundational for any logic, like ours,
    that lives on quantitative bounds.
  \item \href{https://arxiv.org/abs/2607.00034}{\emph{Bayesian
    updates from coalgebraic determinisation}}, Manuel Baltieri and
    Nathaniel Virgo, 2026. Filtering as determinisation --- the kind
    of reframing that makes an old thing suddenly composable. Also a
    datum that the programme's funded orbit reaches Tokyo.
\end{enumerate}

To all of the above: thank you. The best way we know to say it is to
ship things you can check.
